\begin{proofbox}
	\textbf{\textcolor{CProof}{思路提示}}

	利用弦切角定理证明角相等，再由公共角得到三角形相似。

	\medskip
	\textbf{\textcolor{CProof}{正式推导}}
	\begin{description}[style=nextline, leftmargin=2.8em, labelsep=0.5em, itemsep=0.55em]
		\item[\textbf{步骤一（作辅助线）：}] 连接 $AT$、$BT$。
			\par\textbf{理由：}
			构造相似三角形 $\triangle PAT$ 与 $\triangle PTB$。

		\item[\textbf{步骤二（证角等一）：}] 由弦切角定理可得
			\[
				\angle PTA = \angle PBT
			\]
			\par\textbf{理由：}
			弦切角等于所夹弧所对的圆周角。

		\item[\textbf{步骤三（证角等二）：}] 因为 $A$ 在 $P$、$B$ 之间，所以射线 $PA$ 与 $PB$ 重合，故
			\[
				\angle APT = \angle BPT
			\]
			\par\textbf{理由：}
			公共角，$P$、$A$、$B$ 共线且 $A$ 在线段 $PB$ 上。

		\item[\textbf{步骤四（证相似）：}] 在 $\triangle PAT$ 与 $\triangle PTB$ 中，
			\[
				\angle APT = \angle BPT,\quad \angle PTA = \angle PBT
			\]
			所以 $\triangle PAT \sim \triangle PTB$（AA）。
			\par\textbf{理由：}
			两角分别相等，则两三角形相似。

		\item[\textbf{步骤五（得等积）：}] 由相似得对应边成比例：
			\[
				\frac{PA}{PT} = \frac{PT}{PB}
			\]
			整理得
			\[
				PT^{2} = PA\cdot PB
			\]
			\par\textbf{理由：}
			相似三角形对应边成比例，交叉相乘即得等积式。
	\end{description}

	\medskip
	\textbf{\textcolor{CProof}{结论回扣}}

	因此，切割线定理成立：$PT^{2}=PA\cdot PB$，切线长是割线上两段线段的比例中项。
\end{proofbox}
